\chapter*{Úvod}
\phantomsection
\addcontentsline{toc}{chapter}{Úvod}

\begin{comment}

    Úvod studentské práce, např\,\dots
    
    Nečíslovaná kapitola Úvod obsahuje \uv{seznámení} čtenáře s~problematikou práce.
    Typicky se zde uvádí:
    (a) do jaké tematické oblasti práce spadá, (b) co jsou hlavní cíle celé práce a (c) jakým způsobem jich bylo dosaženo.
    Úvod zpravidla nepřesahuje jednu stranu.
    Poslední odstavec Úvodu standardně představuje základní strukturu celého dokumentu.
    
    Tato práce se věnuje oblasti \acs{DSP} (\acl{DSP}), zejména jevům, které nastanou při nedodržení Nyquistovy podmínky pro \ac{symfvz}.%
    \footnote{Tato věta je pouze ukázkou použití příkazů pro sazbu zkratek.}
    
    Šablona je nastavena na \emph{dvoustranný tisk}.
    Nebuďte překvapeni, že ve vzniklém PDF jsou volné stránky.
    Je to proto, aby důležité stránky jako např.\ začátky kapitol začínaly po vytisknutí a svázání vždy na pravé straně.
    %
    Pokud máte nějaký závažný důvod sázet (a~zejména tisknout) jednostranně, nezapomeňte si přepnout volbu \texttt{twoside} na \texttt{oneside}!
\end{comment}

%% Uvod
%kam to spada a jake jsou cile 
V současné době je kladen velký důraz na digitalizaci, automatizaci a efektivní průmyslovou výrobu. K moderním výrobním, logistickým nebo montážním linkám patří manipulační systémy doplněné o funkci strojového vidění. Díky této technologii, jsou manipulátory schopné přesné manipulace s předměty, automatické kontroly kvality, ale hlavně i třídění a trasování výrobku v reálném čase. 
Jedním z cílů semestrální práce je návrh a fyzická realizace vzdáleného ovládání dopravního pásu. Na tento hardwarový základ přímo navazuje stěžejní cíl, kterým je uvedení chytré kamery do provozu, ve vývojovém prostředí Automation Studio. Kamera bude mít za úkol, spolehlivě detekovat, rozpoznat a následně sledovat pohybující se objekty na dopravním páse.


%%DOPLNIT
%% jak bylo dosazeno cilu prace?

%%struktura prace